\chapter{Introduction} \label{Chap1}

\section{Research Objective}

Methane (CH\textsubscript{4}) is a powerful greenhouse gas, with a global
warming potential approximately 28 times that of CO\textsubscript{2} over a
100-year horizon \cite{ipcc2021}. Agriculture is a major anthropogenic source,
with enteric fermentation accounting for a substantial proportion of the
sector's greenhouse gas emissions \cite{fao2023,eea2024}. However, emissions
from managed grasslands vary through time as livestock activity interacts with
soil, meteorological, and management conditions. Understanding this variability
therefore requires prediction of when ecosystem-scale emissions occur and how they 
respond to changes in farm management.

Eddy Covariance (EC) measurements provide a means of observing this integrated
ecosystem-scale signal rather than methane emitted by individual animals alone.
The North Wyke Farm Platform (NWFP) offers a rare record of half-hourly EC
CH\textsubscript{4} flux from managed beef-cattle grasslands, accompanied by
meteorological, soil, hydrological, and livestock-management data
\cite{hawkins2025,cardenas2022}. This creates an opportunity for
data-driven methane prediction, but also a difficult modelling problem: the
record contains extensive gaps, its inputs occur at different temporal
resolutions, and its target is highly variable and dominated by episodic
emission peaks.

Existing machine-learning research on EC CH\textsubscript{4} in a temperate grassland 
setting has primarily addressed gap-filling, in which missing observations are reconstructed within an
observed period \cite{irvin2021,zhu2023a}. Forecasting is a different task:
predictions must be generated sequentially using only information available at
the forecast origin. To the author's knowledge, no previous study has
benchmarked machine-learning forecasting of ecosystem-scale EC
CH\textsubscript{4} flux at a managed temperate grassland. Moreover, forecasting
alone cannot answer counterfactual questions about how emissions might change
under alternative livestock-management or climate conditions. This
``what-if'' capability is also an underdeveloped area of research in agricultural digital twins \cite{purcell2023state,fakeye2024}.

The objective of this dissertation is therefore to determine how effectively
machine-learning models can reconstruct, forecast, and project ecosystem-scale
CH\textsubscript{4} flux at NWFP. It develops a three-stage
pipeline for gap-filling the observed Eddy Covariance record, forecasting future
flux, and projecting modelled responses under varying management and
climate scenarios. Model benchmarking is accompanied by uncertainty
quantification and interpretability so that predictive performance, confidence,
and inferred emission drivers can be evaluated together.

\section{Project Aims and Goals}

\subsection*{Aims:}
Advance research in machine learning for CH\textsubscript{4} flux prediction at 
managed grassland, with a focus on gap-filling, forecasting, and scenario projection
under contrasting livestock-management conditions.

\subsection*{Goals:}
\begin{enumerate}
    \item Generation of a pipeline for gap-filling, forecasting, and scenario projection of CH\textsubscript{4} flux at NWFP.
    \item Benchmarking of multiple model families (from classical statistical baselines, tree ensembles, 
    deep sequence models, to recent tabular foundation models) for gap-filling and forecasting.
    \item Apply uncertainty quantification (quantile machine learning) and interpretability 
    methods (permutation importance) across gap-filling, forecasting, and scenario projection.
    \item Generation of future projection under differing livestock-management practices based on 
    inference of previously generated model and CMIP6 climate projections.
\end{enumerate}

\section{Project Overview and Contributions}

This dissertation contributes three benchmarked results to the field:

\begin{enumerate}
    \item \textbf{A gap-filling benchmark for EC CH\textsubscript{4} flux at a managed
    grassland} (Chapter~\ref{Chap4}), comparing established methods, published baseline
    replications, and newer tabular foundation model families, with uncertainty quantification
    attached. The best configuration reaches an hourly-resolution $R^2$ of 0.429 - 0.686 across all
    three towers, against a published ceiling of $R^2<0.1$ at hourly resolution for the same
    ecosystem and data type \cite{zhu2023a}.

    \item \textbf{A forecasting benchmark for EC CH\textsubscript{4} flux at a managed grassland}
    (Chapter~\ref{Chap5}), the first, to the author's knowledge, applied to this ecosystem--data
    combination at a managed temperate grassland, comparing statistical baselines, tree ensembles, and
    tabular foundation models under identical conditions, with calibrated uncertainty quantification and
    interpretability through permutation importance. The best configuration beats seasonal-mean baseline 
    by a wide margin (MASE$=$0.691), though a dedicated error evaluation shows that spike events remain 
    a challenge for all models tested.

    \item \textbf{Scenario projection under contrasting management and climate conditions for EC CH\textsubscript{4} flux}
    (Chapter~\ref{Chap6}): the forecasting architecture developed in this dissertation is extended
    to produce varying CH\textsubscript{4} projections under livestock and management practices for 
    2 distinct climate scenarios (SSP2-4.5 and SSP5-8.5), based on CMIP6 climate projections.
\end{enumerate}

\section{Report Overview}
The subsequent chapters are structured as follows:
\begin{itemize}
    \item Chapter~\ref{Chap2} (\nameref{Chap2}) reviews the literature establishing the gap that this dissertation 
    aims to address, as well as the North Wyke Farm Platform (NWFP).
    \item Chapter~\ref{Chap3} (\nameref{Chap3}) describes the NWFP dataset, exploratory data analysis, and
    processing required.
    \item Chapter~\ref{ChapModels} (\nameref{ChapModels}) outlines the methodology used for the project, including the 
    models, evaluation metrics, cross-validation procedure, uncertainty quantification and 
    interpretability.
    \item Chapter~\ref{Chap4} (\nameref{Chap4}) presents the results of the gap-filling benchmark; 
    \item Chapter~\ref{Chap5} (\nameref{Chap5}) presents the results of the forecasting benchmark; 
    \item Chapter~\ref{Chap6} (\nameref{Chap6}) presents the results of scenario projection.
    \item Chapter~\ref{Chap7} (\nameref{Chap7}) discusses the principal findings, their limitations and directions for future work before concluding the dissertation.    
    % \item Chapter~\ref{Chap8} (\nameref{Chap8}) concludes with an overview of the entire project.
\end{itemize}
